\section{Introduction}
\label{sec:intro}

A learning system is \emph{recursively self-improving} (RSI) when the thing being improved is also the thing doing the improving: the current model writes its own practice problems, attempts them, keeps the attempts an external verifier confirms, trains on them, and the improved model writes the next round's problems. Taken one step further, the system also revises the \emph{procedure} by which it improves. Recent language-model results in this direction are striking~\cite{zelikman2022star,gulcehre2023rest,chen2024spin,yuan2024selfrewarding,zhao2025absolutezero}, reinforcement learning with verifiable rewards (RLVR) has become the standard recipe for reasoning models~\cite{guo2025deepseekr1,shao2024deepseekmath}, and robot learning has started to follow: a vision-language-action (VLA) policy can be improved by RL in simulation from a sparse success signal alone~\cite{li2025simplevlarl,liu2025rl4vla}.

Two questions are usually left open in such reports. First, \emph{which part of the loop is responsible for the gain?} An RSI loop contains ordinary on-policy RL as a special case; a positive result does not by itself show that recursion, the verifier, or the meta level mattered. Second, \emph{is the gain real?} Benchmarks such as LIBERO~\cite{liu2023libero} have a fixed set of official initial states, and it is tempting to train on them; we find that a widely used VLA-RL checkpoint trained this way is inflated by about 17 points relative to freshly sampled placements of the same tasks (\cref{sec:rq5}).

We treat both questions as experimental-design problems. We build a framework (\cref{sec:method}) in which each ingredient of RSI is an independent switch inside one engine: the verifier can be real or permuted, practice can be generated by the latest policy or by the frozen initial one, the meta level can be absent, a numeric search, or a large language model that reads the round records and edits the training procedure, and the whole budget can be spent in one round of plain RL. We then \emph{preregister} every claim (\cref{sec:setup}): comparisons, bars, statistics and the read-once rule were committed before the runs, every change is a timestamped public amendment, the official test is sealed until all decisions are made, and finals are averaged over replicate evaluations.

\paragraph{Findings.} \Cref{fig:teaser} summarises them.
\begin{enumerate}[nosep,leftmargin=*]
  \item \textbf{The gate makes self-training safe.} Over 31 runs with an OpenVLA-OFT policy on LIBERO-Goal, no gated run ends below its start beyond evaluation noise, while the same recipe collapses to 5.5\% without the gate (\cref{sec:rq1}).
  \item \textbf{The loop's gains are small, and no single ingredient clears its bar.} The loop gains one to two points; recursion contributes $+1.7$ $[-0.3, +3.7]$ and the verifier $+1.2$ $[-0.4, +2.9]$; one round of plain RL at equal compute does as well (\cref{sec:rq2}).
  \item \textbf{A meta level does not help a procedure that already works, and we can show why.} Four variants, including an LLM that may edit nine harness components, are within $\pm 1.2$ points of a fixed procedure; a second registered study on two platforms gives $-0.3$ and $+0.3$. The LLM's diagnoses are correct, but the object loop reaches the same plateau under every procedure in the space (\cref{sec:rq3}).
  \item \textbf{The meta level becomes distinguishable from RL only when the procedure fails.} Starting from a recipe that collapses without the gate, the LLM diagnoses the cause unprompted and its first proposal is kept by the gate on three of four seeds; the registered sealed test of this self-repair study is read on 2026-11-08 (\cref{sec:rq3}).
  \item \textbf{In a second domain the same loop gives larger, significant gains, and recursion pays only at long horizon.} Code world models with four LLMs: verified feedback $+2.6$ to $+5.8$ points; recursion $+2.8$ $[+1.2, +4.5]$ at eight rounds, invisible at four (\cref{sec:rq4}).
  \item \textbf{A published VLA-RL number owes about a fifth of itself to the test states} ($+16.8$ points $[+9.8, +23.4]$; \cref{sec:rq5}).
\end{enumerate}

\paragraph{Contributions.} (i) An RSI framework for VLA policies in which every control is a switch inside one engine, with SFT, behaviour cloning and GRPO as interchangeable update operators and an LLM meta level that edits a declared harness space under an attribution guard. (ii) A preregistered, sealed-test protocol for self-improvement claims, with every amendment public. (iii) The first controlled decomposition of an RSI loop on a VLA benchmark, with its negative results and their mechanism. (iv) A positive recursion result in a second domain at longer horizon. (v) A measurement of train-on-test inflation in a published VLA-RL result. Code, preregistrations and all run records are released.
